
\section{PR-Averaged Error Decomposition and the RR Weight}

We start from the full-window PR weight representation of the RR average.
% The deterministic weight estimates start from the finite-start representation
% of \(\sqrt{n} (\overline{\theta}_n^{(\text{RR}, \alpha)} - \theta^*)\) as a
% noise-weighted sum with a deterministic kernel. The stationary \(n_0 = 0\)
% Berry--Esseen bound proved below is stated for the augmented-chain assembly
% built from the same full-window weights.

\textbf{Depth-one expansion.} Unfolding the error recursion of Chapter 1
gives, for every \(k \ge 1\),
\begin{equation}
\begin{aligned}
\theta_k^{(\alpha)} - \theta^*
  = -\alpha \sum_{l = 1}^k \Gamma_{l + 1 : k}^{(\alpha)} \varepsilon(Z_l)
    + \Gamma_{1 : k}^{(\alpha)} (\theta_0 - \theta^*),
\quad
\Gamma_{l + 1 : k}^{(\alpha)}
  \coloneqq \prod_{j = l + 1}^k (I - \alpha A(Z_j)).
\end{aligned}
\end{equation}
Replacing the random products in the noise sum by their deterministic
counterparts \(B_\alpha^{k - l} \coloneqq (I - \alpha \overline{A})^{k - l}\) and keeping
the initial-product discrepancy separate yields the \textbf{depth-one decomposition}:
\begin{equation}
\begin{aligned}
R_{k, \operatorname{init}}^{(\alpha)}
  \coloneqq \left(\Gamma_{1 : k}^{(\alpha)} - B_\alpha^k\right) (\theta_0 - \theta^*),
\end{aligned}
\label{eq:init-product-remainder}
\end{equation}
\begin{equation}
\begin{aligned}
\theta_k^{(\alpha)} - \theta^*
  = -\alpha \sum_{l = 1}^k B_\alpha^{k - l} \varepsilon(Z_l)
    + B_\alpha^k (\theta_0 - \theta^*)
    + R_{k, \operatorname{init}}^{(\alpha)}
    + R_k^{(\alpha)},
\end{aligned}
\label{eq:depth-one}
\end{equation}
where \(R_k^{(\alpha)} \coloneqq J_k^{(1, \alpha)} + H_k^{(1, \alpha)}\) is the
first noise-driven misadjustment remainder.
% The term \(R_{k, \operatorname{init}}^{(\alpha)}\) vanishes when \(\theta_0 = \theta^*\);
% otherwise it is a finite-start transient and is handled in the burn-in
% transfer theorem, not by the stationary augmented-chain misadjustment bound.
% The leading component \(J_k^{(1, \alpha)}\) has a stationary bias of order
% \(\alpha\), so it is not treated as an \(\alpha^{3 / 2}\) term. The
% Berry--Esseen proof below controls this noise-driven remainder by refining it
% to depth two, where \(J^{(2)} + H^{(2)}\) carries the \(\alpha^{3 / 2}\) moment
% scale. The first sum is the \textbf{depth-zero} term and is the only piece that
% carries the limiting Gaussian.

\textbf{PR averaging produces \(Q_l^{(\alpha)}\).} Recall the PR average
\(\overline{\theta}_n^{(\alpha)} = (n - n_0)^{-1} \sum_{k = n_0}^{n - 1} \theta_k^{(\alpha)}\).
For the stationary result in this chapter, set \(n_0 = 0\). A burned-in average
uses the different deterministic weight
\begin{equation}
\begin{aligned}
Q_{l,n_0}^{(\alpha)}
  = \frac{n}{n - n_0} \alpha \sum_{k = \max(n_0, l)}^{n - 1}
      B_\alpha^{k - l},
\end{aligned}
\end{equation}
when the statistic is normalized as
\(- n^{-1 / 2} \sum_l Q_{l,n_0}^{(\alpha)} \varepsilon(Z_l)\).
% The burned-in non-stationary theorem therefore requires separate weight,
% Poisson, and variance-comparison arguments.
Subtracting \(\theta^*\), substituting the depth-one decomposition above, and
exchanging the order of summation in the depth-zero piece,
\begin{equation}
\begin{aligned}
\sum_{k = 0}^{n - 1} \sum_{l = 1}^k B_\alpha^{k - l} \varepsilon(Z_l)
  = \sum_{l = 1}^{n - 1}
      \left(\sum_{k = l}^{n - 1} B_\alpha^{k - l}\right)
      \varepsilon(Z_l),
\end{aligned}
\end{equation}
gives the deterministic PR weight
\begin{equation}
\begin{aligned}
Q_l^{(\alpha)}
  \coloneqq \alpha \sum_{k = l}^{n - 1} B_\alpha^{k - l}
   = \alpha \sum_{j = 0}^{n - l - 1} B_\alpha^j.
\end{aligned}
\label{eq:Q-definition}
\end{equation}
Multiplying by \(\sqrt{n}\) gives
\begin{equation}
\begin{aligned}
\sqrt{n} \, (\overline{\theta}_n^{(\alpha)} - \theta^*) = W^{(\alpha)} + D^{(\alpha)},
\end{aligned}
\end{equation}
where
\begin{equation}
\begin{aligned}
W^{(\alpha)} \coloneqq -\frac{1}{\sqrt{n}} \sum_{l = 1}^{n - 1} Q_l^{(\alpha)} \, \varepsilon(Z_l),
\end{aligned}
\label{eq:W-alpha}
\end{equation}
and
\begin{equation}
\begin{aligned}
D^{(\alpha)}
  \coloneqq D_{\operatorname{tr}}^{(\alpha)}
     + D_{\operatorname{init}}^{(\alpha)}
     + D_R^{(\alpha)},
\quad \\
D_{\operatorname{tr}}^{(\alpha)}
  \coloneqq \frac{1}{\sqrt{n}} \sum_{k = 0}^{n - 1} B_\alpha^k (\theta_0 - \theta^*),
\quad
D_{\operatorname{init}}^{(\alpha)}
  \coloneqq \frac{1}{\sqrt{n}} \sum_{k = 0}^{n - 1}
      R_{k, \operatorname{init}}^{(\alpha)},
\quad
D_R^{(\alpha)}
  \coloneqq \frac{1}{\sqrt{n}} \sum_{k = 0}^{n - 1} R_k^{(\alpha)}.
\end{aligned}
\end{equation}
The RR combination of the noise-driven part is the misadjustment controlled
below by the Levin depth-two transfer.
% The first sum is the deterministic initial-condition transient. Under the
% full-average convention it is only \(O((\sqrt{n} \alpha a)^{-1})\) in general, so
% it must either be retained explicitly or removed by a centered initialization
% \(\theta_0 = \theta^*\). The second sum is stochastic but still purely
% finite-start: it is controlled after burn-in by random-product stability. The
% third is the source of the leading non-Gaussian correction in the
% Berry--Esseen rate. Its RR-combination
% \(2 D_R^{(\alpha)} - D_R^{(2 \alpha)}\) is the \textbf{misadjustment}
% \(D_1^{\text{mis, RR}}\) controlled below by the Levin depth-two transfer.

\textbf{RR combination produces \(\mathcal{Q}_l^{\text{RR}}\).} By linearity,
\begin{equation}
\begin{aligned}
\sqrt{n} \, (\overline{\theta}_n^{(\text{RR}, \alpha)} - \theta^*)
  = W^{\text{RR}} + D^{\text{RR}},
\end{aligned}
\end{equation}
\begin{equation}
\begin{aligned}
W^{\text{RR}} \coloneqq 2 W^{(\alpha)} - W^{(2 \alpha)}
        = -\frac{1}{\sqrt{n}} \sum_{l = 1}^{n - 1} \mathcal{Q}_l^{\text{RR}} \, \varepsilon(Z_l),
\end{aligned}
\label{eq:W-RR}
\end{equation}
where the \textbf{RR weight} is
\begin{equation}
\begin{aligned}
\mathcal{Q}_l^{\text{RR}} \coloneqq 2 Q_l^{(\alpha)} - Q_l^{(2 \alpha)}.
\end{aligned}
\label{eq:Q-RR-definition}
\end{equation}
% The RR weight is deterministic because both PR averages use the same chain
% realization.
% Since both PR averages share the same noise realization \(\left\{Z_k\right\}\), the bracket
% \(\mathcal{Q}_l^{\text{RR}}\) is a single deterministic matrix kernel: all the stochastic
% content of \(W^{\text{RR}}\) now lives in \(\varepsilon(Z_l)\) alone.

% The two required deterministic controls are
% \(\sum_l \lVert \mathcal{Q}_l^{\text{RR}} - \overline{A}^{-1} \rVert^2\) and
% \(\sum_l \lVert \mathcal{Q}_{l + 1}^{\text{RR}} - \mathcal{Q}_l^{\text{RR}} \rVert\).
% The first controls variance comparison; the second controls the
% Poisson-equation / Abel-summation remainder. Both quantities are estimated in
% Sections 4.3--4.4 below. The closed-form identities of Section 4.2 reduce them
% to elementary operator-norm calculations on \(B_\alpha^m - B_{2 \alpha}^m\).
